% Provenance: copied from the design workspace CambridgePQ/recursion.fast, `report/sections/02-objectives.tex`, on 2026-10-19. Read-only copy; the source is authoritative.
% 02-objectives -- written by W2
\section{Objectives and non-goals}\label{sec:objectives}

This section states what recursion.fast optimizes, which designs it intends to reward, the
principles that govern every later rule, and what is deliberately out of scope.
\Cref{fig:obj-pipeline} shows the pipeline that turns a submission into a score; the rest of the
report specifies each box.

\begin{figure}[!b]
\centering
\begin{tikzpicture}[
  node distance=0.5cm and 0.45cm,
  box/.style={draw,rounded corners=2pt,align=left,font=\scriptsize,inner sep=3pt,text width=3.1cm},
  untr/.style={box,fill=red!4,draw=red!45!black},
  trus/.style={box,fill=blue!4,draw=blue!45!black},
  lean/.style={box,fill=green!6,draw=green!40!black},
  outb/.style={box,fill=black!3,draw=black!50},
  arr/.style={-{Latex[length=1.8mm]},thick},
  lab/.style={font=\scriptsize\itshape,align=center}]
\node[untr] (sub) {\textbf{Submission (PR)}\\
  raw images \code{*.code}, \code{*.data}\\
  \code{claim.json} (counters, shape, scores)\\
  Lean proofs \code{Solution.lean}, \code{Candidate/**}\\
  native indexer + prover (records)\\
  \code{NOTES.md}};
\node[trus, right=of sub] (har) {\textbf{Trusted harness}\\
  export root; source policy;\\
  import allow-list; strict claim parser;\\
  \emph{render} bytes and claim\\
  $\to$ organiser data of the\\ rendered challenge};
\node[lean, right=of har] (chk) {\textbf{Lean checking}\\
  0: \code{.olean} manifest;\\
  1: comparator first, in a\\ fresh tree: statement\\ equality, axioms, two\\ kernel replays;\\
  2: environment gate in a\\ second fresh tree;\\
  manifest after each stage};
\node[outb, below=0.9cm of chk] (score) {\textbf{Score}\\
  the certificate proves the\\
  rendered score equals the\\
  contract's score function;\\
  the scorer takes its root};
\node[trus, below=0.9cm of har] (badge) {\textbf{3: Badge path (unscored)}\\
  own process, after the gate;\\
  emulator on the same bytes;\\
  compiled Lean indexer/prover vs native;\\
  hidden benchmark; \emph{measured} $\le$ \emph{claimed}};
\node[outb, below=0.9cm of sub] (rec) {\textbf{Record}\\
  verified + badge + strict\\
  improvement; contract id;\\
  copied to submissions \code{main}};
\draw[arr] (sub) -- (har);
\draw[arr] (har) -- (chk);
\draw[arr] (chk) -- node[lab,right]{verified} (score);
\draw[arr] (har) -- node[lab,right]{same bytes} (badge);
\draw[arr] (score) -- (badge);
\draw[arr] (badge) -- (rec);
\end{tikzpicture}
\caption{The recursion.fast pipeline. Red: untrusted inputs. Blue: trusted software. Green: Lean
kernel checking. Stages run in the order of Decision~A-3. The proved score never depends on running
candidate code; the badge path runs candidate binaries in a sandbox and gates only record
eligibility. The v0.3 harness implements stages 0--2 and a hook for stage 3; the badge checks and
the emulator are specified but not yet built (\cref{sec:architecture}).}
\label{fig:obj-pipeline}
\end{figure}

\subsection{What is optimized, and why}\label{subsec:obj-what}

\paragraph{Primary objective: the cost of one recursion step.} The brief's \enquote{as efficient as
possible} admits four readings (\cref{sec:brief-review}): the cost of the final verifier, the cost
of one recursion step, proof size, and prover time. recursion.fast optimizes the second. In a
recursive proof system the inner verifier is executed inside a proof at every node of the
recursion tree, so its cost dominates the total prover work, while the final verifier runs once.
The recursion-step cost also subsumes the final-verifier cost, since a decider is itself a
verifier program. Proof size and prover time pull against verifier cost (a lower code rate shrinks
proofs and verifier work and inflates the prover); they enter as constraints and secondary
displays (\cref{sec:cost}), not as an arbitrary exchange rate inside the primary score.

\paragraph{Operationalization.} \enquote{Cost} means a \emph{vector} of instruction counts on one
fixed, contract-owned machine: RV64IM with all W-forms plus the Zbc carry-less multiplications,
with an oracle-query system call that stands for a Merkle opening (\cref{sec:protocol}). The vector
has 18 counters: instructions by class, first-touched memory words and pages, and query counts,
Merkle levels and leaf bytes. A submission proves, in Lean, an upper bound $c_k$ on this vector for
\emph{every accepting run} of its verifier image at each of three size classes
$k\in\{16,20,24\}$ (constraint systems with up to $2^k$ rows and columns). The score weights $c_k$
under four calibrated zkVM cost profiles, adds a static Fiat--Shamir surcharge and a proof-byte
charge, applies a soft penalty on prover-oracle blow-up, and takes the geometric mean over profiles
and classes (\cref{sec:cost}). The score function is a Lean definition of the contract, and the
certificate proves that the claimed score is its value. The score is therefore a theorem about the
exact bytes the deployed verifier would execute, not a measurement. A secondary latency board
scores a deterministic pipeline timing model with fork-join harts (\cref{subsec:cost-lat}).

\paragraph{Why accepting runs.} In recursion the circuit proves \enquote{the verifier accepted
$\pi$}; only accepting executions are ever arithmetized, and a zkVM's trace length is the executed
instruction count. A rejecting run is never proved. The maximum over accepting runs is thus
exactly the quantity a recursion step pays, and a maximum over \emph{all} runs over-approximates it
by the most expensive rejecting path, which no deployment pays except under attack. The bound is
not an average and not a measurement on honest proofs: by perfect completeness it covers every
honest run at every challenge vector, and it also covers every accepting run on an adversarial
transcript, including accepting runs on false statements. A verifier cannot make acceptance cheap
on true statements and expensive on false ones, because it cannot tell them apart; any
\enquote{fast accept} path must itself be sound, and then it is simply a cheaper verifier. Rejecting
runs are bounded by the contract's fuel constant ($2^{36}$ steps per hart), and fuel exhaustion is a
rejection (\cref{subsec:proto-machine}). The price is that the deployed decider's denial-of-service
exposure is bounded only by the fuel: $2^{36}$ steps for an image without harts, and
$(\#\mathrm{spawns}+1)\cdot2^{36}$ with harts, since a child runs on its parent's remaining fuel. An
optional all-runs badge is planned; it is not part of contract v0.3.

\paragraph{Relation to the brief's ambitious variant.} The brief's end state arithmetizes the
verifier and measures a 2-to-1 recursion step. We schedule this as a later score track (T3): wall
time is not reproducible, the RISC-V-to-constraint arithmetizer would be a large trusted artifact,
and the best modelled verifiers (\cref{subsec:bench-trajectory}) are 9--29 times too large for a
self-recursion fixed point inside the $2^{24}$ class (\cref{subsec:rec-t3}). The v1 proxy, the
instruction count of the fold verifier on the contract machine, is exactly what a fixed zkVM would
arithmetize.

\subsection{What counts as an interesting verifier architecture}\label{subsec:obj-interesting}

The competition succeeds if the leaderboard is populated by verifier \emph{architectures}, not by
assembly micro-optimizations of one known design, and not by artifacts of the rules. We call a
submission architecturally interesting if it is
\begin{enumerate}[nosep]
\item \emph{general}: one verifier image for every rank-1 constraint system (R1CS) of the class,
  with the circuit entering only through a short verification key and a committed index oracle
  (holography; \cref{subsec:bg-r1cs});
\item \emph{succinct}: its accepting-run cost is sublinear in the number of nonzero matrix entries,
  which in practice means a polynomial IOP combined with a hash-based polynomial commitment or
  proximity test;
\item \emph{provably sound at the stated level}: round-by-round errors of at most $2^{-128}$ per
  round, proved in Lean in the proven proximity regimes, or in a separate, labelled track under a
  published and named hypothesis (\cref{subsec:sec-hypotheses});
\item \emph{recursion-shaped}: it either verifies a full proof or, in the accumulation track (T2),
  folds two accumulators into one of fixed size and defers work to a decider (\cref{sec:recursion}).
\end{enumerate}
The intended design space is wide: sumcheck-based interactive oracle proofs such as Spartan
\cite{Set20-Spartan} and GKR-style memory checking for the sparse index; Reed--Solomon IOPPs over
binary additive domains such as WHIR \cite{ACFY24-WHIR}, BaseFold \cite{ZCF24-BaseFold},
FRI-Binius \cite{DP24-FRIBinius} and Ligerito \cite{NA25-Ligerito}; tensor-code commitments in the
Ligero/Brakedown line \cite{AHIV17-Ligero,GLSTW23-Brakedown}; holographic compilers in the
Marlin/COS line \cite{CHMMVW20-Marlin,COS20}; and hash-based accumulation such as WARP
\cite{BCFW25-WARP}, Arc \cite{BMNW25-Arc} and Blaze \cite{BCFRRZ25-Blaze}. The cost model is
calibrated so that the main levers are protocol choices. In a modelled Spartan-plus-binary-WHIR
verifier at $2^{24}$, the openings of the index commitment and of the witness commitment take
93\% of the cost; that cost is set by the proven query count, the leaf size and folding schedule,
and the size of the index commitment. Hand-tuning the extension-field kernels moves the total by
about 15\% (\cref{sec:cost}).

Equally explicit is what the rules must \emph{not} reward: special-casing the benchmark circuit or
its instances; lookup tables in an unbounded image or verification key; free computation by the
trusted harness (for example, powers of a challenge precomputed on the input tape); typed prover
messages that silently promise a property, such as low degree, that the deployed verifier never
checks; and protocols whose per-round error is small only because the number of rounds is
unbounded. Each of these is a concrete attack found during design review, and each is closed by a
rule cited where it is introduced.

\subsection{Design principles}\label{subsec:obj-principles}

\begin{description}[leftmargin=0pt,itemsep=2pt]
\item[P1. A small, auditable trusted base.] Everything a reader must accept is listed in one table
  (\cref{tab:arch-trusted}). The trusted Lean contract is 796 code lines (v0.3; v0.2 had 582) over
  Lean core, Mathlib and a vendored RISC-V model (listings in \cref{sec:appendix-contract}),
  and imports neither ArkLib nor CompPoly; bridges to those libraries are kernel-checked but
  untrusted. The growth from v0.2 is the cost model and the latency board: counters, timer, harts,
  profiles and score functions. The security definitions are statements about finite sets
  (\enquote{at most an $\eps$-fraction of challenges is bad}) with an existentially quantified
  state function, and no probability monad in the audited text.
\item[P2. A large design space.] The contract fixes the relation, the fields, the machine and the
  shapes of messages and challenges, and nothing about the protocol. Any public-coin holographic
  interactive oracle proof expressible with raw field-vector messages and contract-typed challenges
  is admissible. Organizers ship an untrusted, kernel-checked library (machine facts, a symbolic
  executor, verified field routines, IOP components, eventually a commitment scheme) so that the
  proof burden does not select designs.
\item[P3. Scores are proved, not measured.] The score is a Lean function of scalars that appear
  in kernel-checked theorem statements, and its value is itself a certificate clause. Measurements on hidden benchmarks are unscored
  demonstrations, a differential test of the trusted machine semantics, and a gate on record
  eligibility (\cref{subsec:arch-badge}).
\item[P4. A class-uniform cost theorem.] Knowledge soundness does not forbid
  hard-coding (the extractor is existential and may carry a benchmark witness). Overfitting is
  therefore excluded by the \emph{cost} statement, which quantifies over every index of a class
  shape, every public input, every prover message and every challenge, so that a fast path for one
  circuit cannot lower the score (\cref{subsec:bench-hardcode}).
\item[P5. Honest labelling.] Every record states precisely what is proved. With Lean extractors
  of unbounded running time, the base-level knowledge certificate is equivalent to round-by-round
  \emph{soundness}, and records say so (\cref{subsec:sec-extractor}); records under a named
  hypothesis are flagged with its version; one-step accumulation results are labelled one-step until
  the tree theorem exists; the gap from the interactive model to the deployed hash-based argument is
  stated as cited theorems and named deployment assumptions, never as Lean hypotheses.
\end{description}

\subsection{Non-goals}\label{subsec:obj-nongoals}

\begin{description}[leftmargin=0pt,itemsep=2pt]
\item[Zero knowledge.] Recursion needs succinctness and knowledge soundness, not privacy; zero
  knowledge would add a second simulation-based definition to the trusted text and is out of scope
  in v1.
\item[Quantum adversaries.] The statements proved in Lean are information-theoretic: they hold
  against unbounded adversaries. The compilation to a hash-based argument is cited in
  the classical random-oracle model only; a quantum random-oracle analysis \cite{CMS19} has worse
  constants and is not claimed.
\item[A prover-time score in v1.] The honest prover is a Lean specification proved complete; its
  running time is not observable from a Lean term. Prover cost is constrained statically
  (oracle-length caps and a soft blow-up factor in the score) and practically (a native prover must
  produce accepting proofs within a wall-clock and memory budget before a record is granted). A
  proved prover-cost bound would require a RISC-V prover with a refinement proof, which is far beyond
  any check budget (\cref{sec:cost}).
\item[Arithmetized recursion in v1.] Proving \enquote{the verifier accepted} inside the next proof
  needs a trusted RISC-V-to-constraint compiler and, for hash-compiled verifiers, a recursion
  heuristic that no random-oracle theorem justifies \cite{BCG24-RelROM}. v1 scores the fold
  verifier's cost on the machine; full recursion is a target schema (security level L4, score track
  T3; \cref{sec:recursion}).
\end{description}
